\chapter{Gamma Scalping and Event Volatility}\label{ch:s2:gamma-scalping-and-event-volatility}

Before earnings, a stock's options price a move, and then the stock moves more or less than that. Gao, Xing and Zhang found that at-the-money
straddles on individual stocks generally lose money, yet the same straddles bought three days before an earnings announcement and held to the
announcement date earned a highly significant 3.34\% on average. Owning a straddle and hedging it pays the difference between the move priced and the
move delivered, less the cost of hedging. This chapter plants earnings gaps in synthetic stocks whose implied event variance is too low by a
controlled amount. On them the \omterm{def:s2:gamma-scalping-and-event-volatility:straddle}{event straddle} earns 2.1\% of its premium per event, which is a Sharpe ratio of about 1 for a book of 200 events a year.
One event tells you almost nothing, though: the return's standard deviation is 29\% of the premium. The hedging frequency that does best after costs
is once a day. The build is \texttt{firm.eventvol}.

\section{Gamma scalping as a bet on realised against implied}

A long option hedged in delta earns, over a short interval, about $\tfrac12\Gamma S^2(\sigma_r^2 - \sigma_i^2)\,\mathrm{d}t$: gamma times the
difference between the variance delivered and the variance the option was priced at (Book~5, chapter~25). Each hedge sells the stock after it has
risen and buys it after it has fallen. The trades bank the path's variance, and theta pays for the implied variance. The trade is a bet on realised
against implied volatility, and the delta hedge is how the bet is collected.

The chapter's synthetic stocks have a diffusive vol around 30\%, and their implied variance carries a premium of 25\% over it. Buying a 21-day
at-the-money straddle and hedging it for 20 trading days, with no event, loses the premium, whatever the hedging:

\begin{center}
\small
\begin{tabular}{@{}lccc@{}}
\toprule
2\,000 straddles, no event & mean return (\% of premium) & sd (\%) & hedging cost (\%)\\
\midrule
never hedged & $-12.1$ & 62.4 & 0.0\\
daily & $-12.4$ & 15.3 & 0.7\\
twice a day & $-13.0$ & 11.5 & 1.0\\
every two hours & $-13.4$ & 8.5 & 1.5\\
hourly & $-14.1$ & 6.7 & 2.1\\
when delta moves 0.05 & $-13.7$ & 7.2 & 1.6\\
when delta moves 0.10 & $-13.0$ & 8.6 & 1.2\\
when delta moves 0.20 & $-12.6$ & 12.6 & 0.7\\
\bottomrule
\end{tabular}
\end{center}

Hedging does not change what the trade expects to earn. That is set by implied against realised variance. What hedging changes is how noisily the
trade earns it. Unhedged, the straddle is a directional bet with a standard deviation of 62\% of its premium. Hedged hourly, the standard deviation is
6.7\%, and 2.1\% of the premium goes to the stock trades that hedge it. The hedges cost two basis points of the stock traded, an assumption of this
chapter.

A band rule hedges only when the delta has moved by more than a threshold since the last hedge, so it trades when the path moves rather than when the
clock strikes (\cref{fig:s2:gamma-scalping-and-event-volatility:scalping}). At similar costs it gets a lower standard deviation than the clock: a band of
0.10 costs 1.2\% for 8.6\%, where hedging every two hours costs 1.5\% for 8.5\%.

\begin{omfigure}
\begin{tikzpicture}
\begin{axis}[width=10cm,height=5.6cm,xlabel={\small hedging cost (\% of premium)},ylabel={\small sd of return (\% of premium)},
  tick label style={font=\footnotesize},xmin=0.2,xmax=2.4,ymin=5,ymax=17,grid=major,grid style={gray!25},table/col sep=comma,
  nodes near coords,point meta=explicit symbolic,every node near coord/.append style={font=\scriptsize}]
\addplot[only marks,mark=*,omThm,every node near coord/.append style={anchor=west}] table[x=cost,y=sd,meta=rule]{figdata/strategies-2/04-gamma-scalping-and-event-volatility/scalping_clock.csv};
\addplot[only marks,mark=square*,omDef,every node near coord/.append style={anchor=east}] table[x=cost,y=sd,meta=rule]{figdata/strategies-2/04-gamma-scalping-and-event-volatility/scalping_band.csv};
\end{axis}
\end{tikzpicture}

\omcaption{Hedging a long 21-day straddle for 20 days on the synthetic stocks, with no event: each rule's hedging cost against the standard deviation
of the return, both as a share of the premium. Clock rules (circles, daily to hourly) and
band rules (squares, delta moves of 0.05 to 0.20). Data:
\texttt{s2\_eventvol.hedging(False)}.}\label{fig:s2:gamma-scalping-and-event-volatility:scalping}
\end{omfigure}

\section{Implied moves around events}

An earnings announcement comes on a known day and adds a lump of variance on that day. An expiry after the event carries it; an expiry before does not.
For an expiry $\tau$ after the event, with diffusive implied vol $\sigma_d$ and event variance $E$,
\[
\sigma^2(\tau)\,\tau = \sigma_d^2\,\tau + E,
\]
so the event variance lifts short expiries' vols more than long ones', and two expiries after the event are enough to solve for both $\sigma_d$ and
$E$. Dubinsky, Johannes, Kaeck and Seeger built reduced-form models and estimators on this idea to separate the uncertainty about earnings
announcements from day-to-day volatility. The anticipated uncertainty was quantitatively large, varied over time and was informative about future
volatility.

\begin{definition}[Implied move]\label{def:s2:gamma-scalping-and-event-volatility:move}
The \emph{implied move}\index{implied move} of a scheduled event is the expected absolute return on the event day that option prices imply. Under a
normal event return with the event variance $E$ extracted from the term structure it is $\sqrt{2/\pi}\,\sqrt{E}$. Traders often approximate it by
the price of the at-the-money straddle expiring just after the event, as a share of the stock price.
\end{definition}

On one of the chapter's stocks, with a diffusive implied variance and an implied event variance of 0.00135, a 5-day expiry trades at 39.5\% and a
21-day expiry at 32.3\%. The two vols give the event variance back exactly (\texttt{event\_variance}): an event standard deviation of 3.67\%.

\begin{definition}[Event straddle]\label{def:s2:gamma-scalping-and-event-volatility:straddle}
An \emph{event straddle}\index{event straddle} is a long at-the-money straddle bought shortly before a scheduled event and sold soon after it, usually
delta-hedged, which profits when the move delivered on the event exceeds the \omterm{def:s2:gamma-scalping-and-event-volatility:move}{implied move} by more than the time decay and costs of holding it.
\end{definition}

The synthetic market plants the mispricing. Each event's true move has a standard deviation around 5\%, drawn event by event. The market prices an
implied event variance equal to the true one times 0.70, times a lognormal error: the \omterm{def:s2:gamma-scalping-and-event-volatility:move}{implied move} is too low on average, and wrong by a different
amount each time. Across the 2\,000 events the average absolute gap was 4.35\% and the average \omterm{def:s2:gamma-scalping-and-event-volatility:move}{implied move} 3.62\%. The median \omterm{def:s2:gamma-scalping-and-event-volatility:move}{implied move} was 0.81 of
the true one. The book buys the 21-day straddle at the close three days before the event, hedges it by a rule, and sells it at the event day's close.
Options cost 2\% of the premium for the round trip, and hedges two basis points.

\begin{center}
\small
\begin{tabular}{@{}lccccc@{}}
\toprule
2\,000 events & mean (\%) & sd (\%) & mean/sd & hedging cost (\%) & 200 events a year\\
\midrule
never hedged & 2.0 & 31 & 0.064 & 0.00 & 0.91\\
daily & 2.1 & 29 & 0.073 & 0.06 & 1.04\\
twice a day & 2.0 & 28 & 0.071 & 0.18 & 1.00\\
every two hours & 1.8 & 27 & 0.064 & 0.23 & 0.91\\
hourly & 1.7 & 27 & 0.063 & 0.30 & 0.89\\
when delta moves 0.05 & 1.7 & 27 & 0.063 & 0.24 & 0.89\\
when delta moves 0.10 & 2.0 & 28 & 0.070 & 0.18 & 1.00\\
when delta moves 0.20 & 1.9 & 28 & 0.068 & 0.13 & 0.96\\
\bottomrule
\end{tabular}
\end{center}

The last column scales the per-event ratio by $\sqrt{200}$, a book of fifty names with four reports each, and assumes the events are independent,
which they are in this simulation and are not in a market. The gap arrives at the open, before anyone can hedge, so no hedging rule captures it. Hedging
only trims the noise of the days around it. Daily hedging does best after costs: hedging more often cuts the standard deviation from 29\% to 27\% and
costs more than it saves.

Planting the effect matters more than any rule. With the implied event variance right on average, the daily-hedged book loses 2.2\% per event. Priced
15\% too low, it loses 0.1\%. Only at 30\% too low does it earn its 2.1\%. The straddle has to beat its own theta, the diffusive premium and its costs
before the event pays anything.

\section{The volatility crush after events}

\begin{definition}[Volatility crush]\label{def:s2:gamma-scalping-and-event-volatility:crush}
The \emph{volatility crush}\index{volatility crush} is the fall in implied volatility after a scheduled event, when the event's variance leaves the
options' remaining life; it is largest for the shortest expiries, whose vol the event had lifted most.
\end{definition}

On the synthetic stocks the median 21-day implied vol was 36.6\% when the straddle was bought and 32.0\% after the event, a median fall of 8.6\%. The
straddle holder loses that vega on the event day whatever the stock does, so the gap must pay for it.

The same pattern shows on the index around scheduled Federal Reserve statements. On the 125 scheduled FOMC statement days from January 2011 to
September 2026, the S\&P 500's mean absolute log return was 0.83\%, against 0.71\% on other days. Its mean squared return was 1.146 times larger,
an extra event move with a standard deviation of about 0.41\%. Cboe's 9-day VIX fell on 65.6\% of statement days, against 53.1\% of other days, by
3.05\% on average (log change), against a rise of 0.09\% on other days. The 30-day VIX, which the statement lifts less, fell by 1.10\% on average.
A macro event is small next to an earnings report, and the crush it leaves in the shortest vol is still visible.

\section{Hedging frequency and path}

The \omterm{def:s2:gamma-scalping-and-event-volatility:move}{implied move}'s error decides each event (\cref{fig:s2:gamma-scalping-and-event-volatility:error}). Sorted into fifths by the error, the log of the
\omterm{def:s2:gamma-scalping-and-event-volatility:move}{implied move} over the true one, the daily-hedged straddle earned 8.0\% of its premium when the \omterm{def:s2:gamma-scalping-and-event-volatility:move}{implied move} was 0.46 too low in logs. Where the \omterm{def:s2:gamma-scalping-and-event-volatility:move}{implied
move} was about right (an error of 0.05), it lost 4.5\%. The correlation between the error and the return, event by event, is only $-0.13$, because the
move delivered is a single draw: a well-priced event can still gap three standard deviations.

\begin{omfigure}
\begin{tikzpicture}
\begin{axis}[width=10cm,height=5.4cm,ybar,bar width=16pt,xlabel={\small fifth of events by the implied move's error (1: most underpriced)},
  ylabel={\small mean return (\% of premium)},tick label style={font=\footnotesize},xmin=0.4,xmax=5.6,ymin=-6,ymax=10,xtick={1,2,3,4,5},
  grid=major,grid style={gray!25},table/col sep=comma]
\addplot[fill=omThm!70,draw=omThm] table[x=quintile,y=mean]{figdata/strategies-2/04-gamma-scalping-and-event-volatility/error.csv};
\end{axis}
\end{tikzpicture}

\omcaption{The daily-hedged \omterm{def:s2:gamma-scalping-and-event-volatility:straddle}{event straddle}'s mean return by fifth of the 2\,000 synthetic events, sorted by the \omterm{def:s2:gamma-scalping-and-event-volatility:move}{implied move}'s error (log of implied
over true move; mean errors $-0.46$, $-0.30$, $-0.21$, $-0.11$ and 0.05). Data: \texttt{s2\_eventvol.by\_error}.}\label{fig:s2:gamma-scalping-and-event-volatility:error}
\end{omfigure}

The path matters as well as the frequency. A clock rule hedges quiet days and busy days alike, while a band rule waits for the delta to move. Around
an event the path has a structure: the diffusive days before, then the gap at the open. The gap cannot be hedged, so any rule that hedges the days
before it is paying to trim variance it could have left alone. That is why the daily rule and the 0.10 band tie, and why hourly hedging loses 0.4\%
of the premium to the daily rule over three days.

\section{Strategy files}

\begin{strategyfile}{Earnings event straddle}\label{strat:s2:gamma-scalping-and-event-volatility:earnings}
\sfield{Who pays you, and why} Option sellers who underestimate the uncertainty of an announcement, more so when a firm's signals are noisy and its
options costly to trade.
\sfield{Instruments and venues} Listed single-stock options, at-the-money straddles with the first expiry after the report.
\sfield{Signal} \omterm{def:s2:gamma-scalping-and-event-volatility:move}{Implied move} against a forecast of the move: past surprises, kurtosis, size, dispersion of analysts.
\sfield{Sizing and execution} Many small positions across names and reporting seasons; bought a few days ahead, sold after the report.
\sfield{Costs} Option spreads, widest where the effect is largest.
\sfield{How it dies} Crowding that raises \omterm{def:s2:gamma-scalping-and-event-volatility:move}{implied moves}; costs that absorb the premium.
\sfield{Horizon, capacity, infrastructure} Days per trade; an earnings calendar and option surfaces for hundreds of names.
\sfield{Backtest honestly} Quoted spreads, not mid prices; announcement dates as known then; the crush in the exit price.
\sfield{Sources} Gao, Xing and Zhang (2018): 3.34\% on average from three days before to the announcement date; Dubinsky, Johannes, Kaeck and Seeger
(2019); this chapter: 2.1\% per event, about 1.0 for 200 events a year.
\end{strategyfile}

\begin{strategyfile}{Post-event volatility sale}\label{strat:s2:gamma-scalping-and-event-volatility:postevent}
\sfield{Who pays you, and why} Holders of event options who sell after the event, and the diffusive variance premium that remains.
\sfield{Instruments and venues} Single-stock or index options after the report or the statement.
\sfield{Signal} Remaining implied vol against the forecast diffusive vol once the event is past.
\sfield{Sizing and execution} Delta-hedged short straddles or variance, sized by a gap stress.
\sfield{Costs} Option spreads; hedging.
\sfield{How it dies} Follow-on news (guidance, restatements) that brings a second event.
\sfield{Horizon, capacity, infrastructure} Days to weeks; surfaces updated after each event.
\sfield{Backtest honestly} Exit and entry at quoted prices after the event, not at the close before.
\sfield{Sources} No performance figure verified.
\end{strategyfile}

\begin{strategyfile}{Gamma scalping with hedge bands}\label{strat:s2:gamma-scalping-and-event-volatility:bands}
\sfield{Who pays you, and why} Option sellers who price vol below what the path delivers.
\sfield{Instruments and venues} Long options on liquid underlyings; the stock or future for hedges.
\sfield{Signal} Forecast realised against implied vol.
\sfield{Sizing and execution} Delta bands set by cost against variance; hedging on moves, not on the clock.
\sfield{Costs} Stock or futures trading at every hedge.
\sfield{How it dies} Implied above realised on average (the variance premium): a long gamma book loses it.
\sfield{Horizon, capacity, infrastructure} Days to weeks; a hedging engine with live deltas.
\sfield{Backtest honestly} Intraday paths; hedging costs per trade.
\sfield{Sources} This chapter: without an event, the long straddle loses 12\% to 14\% of its premium whatever the rule; hedging sets the noise.
\end{strategyfile}

\begin{strategyfile}{Macro-event variance}\label{strat:s2:gamma-scalping-and-event-volatility:macro}
\sfield{Who pays you, and why} Buyers of protection over central-bank decisions and data releases, and the crush afterwards.
\sfield{Instruments and venues} Short-dated index options and weekly expiries around statement days.
\sfield{Signal} The event variance in the short expiries against the index's past moves on such days.
\sfield{Sizing and execution} Small; sized by the largest past event-day moves.
\sfield{Costs} Short-dated index option spreads.
\sfield{How it dies} A surprise decision; an unscheduled meeting.
\sfield{Horizon, capacity, infrastructure} Days; a calendar of scheduled events.
\sfield{Backtest honestly} Scheduled dates as published then; unscheduled meetings excluded from the signal and kept in the risk.
\sfield{Sources} Federal Reserve calendars and Cboe histories (this chapter: 125 statement days, VIX9D down 3.05\% on average).
\end{strategyfile}

\section{Tutorial: the move priced and the move delivered}\label{tut:s2:gamma-scalping-and-event-volatility:tut}

\begin{tutorial}
\textbf{Goal.} Simulate event windows, price and hedge \omterm{def:s2:gamma-scalping-and-event-volatility:straddle}{event straddles} by several rules, and measure the crush and the FOMC-day statistics.
\textbf{End state:} the two tables and the two figures.
\begin{tutsteps}
\item \textbf{The straddle and its hedges}.
\omcode{code/firm/eventvol/firm_eventvol.py}{87}{118}{A long straddle per window, delta-hedged by clock or band, with costs.}\label{lst:s2:gamma-scalping-and-event-volatility:pnl}
\item \textbf{The hedging rules}.
\omcode{code/strategies-2/04-gamma-scalping-and-event-volatility/python/s2_eventvol.py}{34}{51}{Each rule's return, noise and cost.}\label{lst:s2:gamma-scalping-and-event-volatility:rules}
\item \textbf{Run} \texttt{s2\_fetch\_fomc.py} once, then \texttt{hedging()}, \texttt{hedging(False)}, \texttt{by\_error()}, \texttt{crush()} and
\texttt{fig\_eventvol.py}.
\end{tutsteps}
\textbf{What to change next.} Let the \omterm{def:s2:gamma-scalping-and-event-volatility:move}{implied move}'s error depend on something observable (past surprises) and trade only the fifth most likely
underpriced; make events correlated across names in a season; add a gap at the open on non-event days.
\end{tutorial}

\section{Build: event volatility}\label{bld:s2:gamma-scalping-and-event-volatility:eventvol}

\begin{build}
\bfield{Purpose} Event variance from the term structure, \omterm{def:s2:gamma-scalping-and-event-volatility:move}{implied moves}, event windows with gaps at the open, and discretely hedged straddles.
\bfield{Interface} \texttt{EventConfig(\dots)}, \texttt{event\_variance(v1, t1, v2, t2)}, \texttt{implied\_move(event\_var)}, \texttt{windows(cfg,
days, event)}, \texttt{straddle\_pnl(w, cfg, every, band)}.
\bfield{Rules} The gap happens at the event day's open, between two hedges; the event variance leaves the implied vol after the gap; costs on every
hedge and on the options.
\bfield{Acceptance tests} \texttt{code/firm/eventvol/tests/}: the event variance recovered exactly from two expiries; the crush after every event and
none without one; on a still path, only time decay and costs.
\bfield{Stretch} Several events per window; jumps on non-event days; stochastic diffusive vol.
\end{build}

\begin{omsources}
\item A.~Dubinsky, M.~Johannes, A.~Kaeck and N.~J.~Seeger, ``Option pricing of earnings announcement risks'', \emph{Review of Financial Studies}
32(2), 2019.
\item C.~Gao, Y.~Xing and X.~Zhang, ``Anticipating uncertainty: straddles around earnings announcements'', \emph{Journal of Financial and Quantitative
Analysis} 53(6), 2018.
\item Federal Reserve, FOMC meeting calendars, 2011--2026.
\item Cboe Global Markets, daily histories of SPX, VIX and VIX9D.
\end{omsources}

\section{Exercises}

\begin{exercise}[$\star$]\label{exo:s2:gamma-scalping-and-event-volatility:1}
An event's variance is 0.0025. What is its \omterm{def:s2:gamma-scalping-and-event-volatility:move}{implied move}?
\end{exercise}

\begin{exercise}[$\star$]\label{exo:s2:gamma-scalping-and-event-volatility:2}
A stock's diffusive implied vol is 30\% and an event adds variance 0.0025. What is the implied vol of an expiry 10 trading days away, after the event?
\end{exercise}

\begin{exercise}[$\star$]\label{exo:s2:gamma-scalping-and-event-volatility:3}
Why can no hedging rule capture an earnings gap?
\end{exercise}

\begin{exercise}[$\star\star$]\label{exo:s2:gamma-scalping-and-event-volatility:4}
Gao, Xing and Zhang found that straddles generally lose but \omterm{def:s2:gamma-scalping-and-event-volatility:straddle}{event straddles} gain. Reconcile the two.
\end{exercise}

\begin{exercise}[$\star\star$]\label{exo:s2:gamma-scalping-and-event-volatility:5}
Why does a band rule get less noise for the same cost than a clock rule?
\end{exercise}

\begin{exercise}[$\star\star$]\label{exo:s2:gamma-scalping-and-event-volatility:6}
The S\&P 500's mean squared return on FOMC days is 1.146 times that of other days. Why is the macro-event trade small next to the earnings trade?
\end{exercise}

\begin{exercise}[$\star\star\star$]\label{exo:s2:gamma-scalping-and-event-volatility:7}
\emph{Coding.} Run \texttt{hedging(True, 0.0)} and \texttt{hedging(True, -0.15)}. How mispriced must the event be for the daily-hedged straddle to
break even, and why is the answer not zero?
\end{exercise}

\begin{exercise}[$\star\star\star$]\label{exo:s2:gamma-scalping-and-event-volatility:8}
\emph{Find the flaw.} ``Our \omterm{def:s2:gamma-scalping-and-event-volatility:straddle}{event straddles} earn 2.1\% per event with a Sharpe ratio of 1.04 a year, so we can size the book on that Sharpe ratio.''
\end{exercise}

\section{Problem: The Move Priced and the Move Delivered}

\begin{problem}[{Weekend problem --- trading events}]\label{pb:s2:gamma-scalping-and-event-volatility:1}
The chapter's synthetic events, the FOMC record and the public research.

\medskip\noindent\textbf{Part I --- Gamma.}
\begin{enumerate}
\item Write the hedged option's P\&L over a short interval and explain it.
\item Why does hedging not change what a long straddle expects to earn?
\item What did the no-event straddle earn under each hedging rule?
\item Define a band rule and compare it with a clock rule.
\end{enumerate}

\medskip\noindent\textbf{Part II --- Events.}
\begin{enumerate}[resume]
\item How is the event variance extracted from the term structure?
\item Define the \omterm{def:s2:gamma-scalping-and-event-volatility:move}{implied move} and the \omterm{def:s2:gamma-scalping-and-event-volatility:straddle}{event straddle}.
\item What did Dubinsky and co-authors, and Gao, Xing and Zhang, find?
\item How does the synthetic market misprice events?
\end{enumerate}

\medskip\noindent\textbf{Part III --- The crush.}
\begin{enumerate}[resume]
\item Define the \omterm{def:s2:gamma-scalping-and-event-volatility:crush}{volatility crush} and give its size on the synthetic stocks.
\item What happens to the S\&P 500 and to VIX9D on FOMC days?
\item Why does the 30-day VIX fall less than the 9-day?
\item What does the crush cost the straddle holder?
\end{enumerate}

\medskip\noindent\textbf{Part IV --- The verdict.}
\begin{enumerate}[resume]
\item State the \emph{named result}: the \omterm{def:s2:gamma-scalping-and-event-volatility:straddle}{event straddle}'s return per event against the \omterm{def:s2:gamma-scalping-and-event-volatility:move}{implied move}'s error, and the hedging frequency that
maximises it after costs.
\item Give the return by fifth of the \omterm{def:s2:gamma-scalping-and-event-volatility:move}{implied move}'s error.
\item Why is the event-by-event correlation only $-0.13$?
\item How mispriced must events be for the trade to pay?
\item What does the 200-events-a-year Sharpe ratio assume?
\item How would you backtest an earnings straddle book honestly?
\item Which strategy file is closest to the variance premium of chapter~1?
\item In one sentence: what does an \omterm{def:s2:gamma-scalping-and-event-volatility:straddle}{event straddle} buy?
\end{enumerate}
\end{problem}

\section{Interview questions}

\begin{interviewq}[$\star$ \iqroles{trader}]\label{iq:s2:gamma-scalping-and-event-volatility:1}
What is gamma scalping, and when does it make money?
\end{interviewq}

\begin{interviewq}[$\star\star$ \iqroles{trader}]\label{iq:s2:gamma-scalping-and-event-volatility:2}
A stock reports tomorrow. How do you read the \omterm{def:s2:gamma-scalping-and-event-volatility:move}{implied move} from the option chain?
\end{interviewq}

\begin{interviewq}[$\star\star$ \iqroles{researcher}]\label{iq:s2:gamma-scalping-and-event-volatility:3}
How would you test whether earnings moves are underpriced?
\end{interviewq}

\begin{interviewq}[$\star\star$ \iqroles{risk}]\label{iq:s2:gamma-scalping-and-event-volatility:4}
A book is short \omterm{def:s2:gamma-scalping-and-event-volatility:straddle}{event straddles} across 200 names in earnings season. What is its risk, and how would you stress it?
\end{interviewq}

\begin{interviewq}[$\star\star$ \iqroles{developer}]\label{iq:s2:gamma-scalping-and-event-volatility:5}
Design a delta-hedging engine with band rules. What must it know, and what can go wrong at the open?
\end{interviewq}

\begin{interviewq}[$\star\star\star$ \iqroles{researcher}]\label{iq:s2:gamma-scalping-and-event-volatility:6}
Two expiries after an event, $\tau_1 < \tau_2$, have implied vols $\sigma_1$ and $\sigma_2$. Derive the event variance and the diffusive vol, and
say what goes wrong if an expiry falls before the event.
\end{interviewq}
